\documentclass[12pt,a4paper]{article}
\usepackage[utf8]{inputenc}
\usepackage[indonesian]{babel}
\usepackage{geometry}
\usepackage{graphicx}
\usepackage{float}
\usepackage{listings}
\usepackage{xcolor}
\usepackage{amsmath}
\usepackage{booktabs}
\usepackage{hyperref}
\usepackage{fancyhdr}

\geometry{a4paper, left=3cm, right=3cm, top=3cm, bottom=3cm}

% Setup untuk kode Python
\lstset{
    language=Python,
    basicstyle=\ttfamily\small,
    keywordstyle=\color{blue},
    commentstyle=\color{green!60!black},
    stringstyle=\color{red},
    numbers=left,
    numberstyle=\tiny\color{gray},
    stepnumber=1,
    numbersep=5pt,
    backgroundcolor=\color{gray!10},
    showspaces=false,
    showstringspaces=false,
    showtabs=false,
    frame=single,
    tabsize=2,
    captionpos=b,
    breaklines=true,
    breakatwhitespace=false,
    escapeinside={\%*}{*)},
    xleftmargin=2em,
    framexleftmargin=1.5em
}

% Header dan Footer
\pagestyle{fancy}
\fancyhf{}
\fancyhead[L]{Laporan Praktikum Mandiri ANN}
\fancyhead[R]{Machine Learning}
\fancyfoot[C]{\thepage}

\begin{document}

\noindent
\textbf{Praktikum Mandiri - MNIST Handwritten Digits} \hfill \textbf{Raffa Yuda Pratama / 0110224081}

\vspace{1em}

\begin{center}
\textbf{\large Implementasi Multi-Layer Perceptron (MLP)} \\
\textbf{\large untuk Klasifikasi Dataset Citra MNIST Handwritten Digits}
\end{center}

\vspace{1em}

\section{Pendahuluan}

Artificial Neural Network (ANN) dengan arsitektur Multi-Layer Perceptron (MLP) merupakan salah satu algoritma deep learning yang efektif untuk klasifikasi multi-kelas. Praktikum mandiri ini mengimplementasikan MLP pada dataset citra MNIST yang berisi gambar tulisan tangan angka 0-9. Tujuan praktikum ini adalah untuk memperluas pemahaman mahasiswa tentang konsep ANN dengan mengaplikasikannya pada data citra (image dataset) yang berbeda karakteristiknya dari dataset tabular seperti Titanic.

\section{Dataset}

Dataset MNIST Handwritten Digits memiliki karakteristik:
\begin{itemize}
    \item Sumber: Kaggle (\texttt{handwritten-digits-0-9})
    \item Total sampel: 5000 gambar (dibatasi 500 gambar per kelas untuk efisiensi)
    \item Dimensi gambar: 28x28 pixels, grayscale
    \item Jumlah kelas: 10 kelas (digit 0-9)
    \item Format: Image files dalam folder terpisah per digit
    \item Target: Multi-class classification (10 kelas)
\end{itemize}

Dataset ini berbeda dengan dataset tabular karena setiap data point adalah gambar 2D yang harus di-flatten menjadi vektor 1D sebelum diproses oleh MLP.

\section{Implementasi}

\subsection{Download Dataset dari Kaggle}

\begin{lstlisting}[caption=Konfigurasi Kaggle API]
kaggle = '../../Data/kaggle.json'

!mkdir -p ~/.kaggle
!cp {kaggle} ~/.kaggle/
!chmod 600 ~/.kaggle/kaggle.json
\end{lstlisting}

\textbf{Penjelasan:} Setup Kaggle API credentials dengan menyalin file \texttt{kaggle.json} ke direktori \texttt{~/.kaggle/} dan mengatur permission file agar aman.

\begin{lstlisting}[caption=Download dan Ekstraksi Dataset]
!kaggle datasets download -d olafkrastovski/handwritten-digits-0-9

from zipfile import ZipFile
import os

file_name = 'handwritten-digits-0-9.zip'
extract_path = 'handwritten_digits'
os.makedirs(extract_path, exist_ok=True)

with ZipFile(file_name, 'r') as zip_ref:
    zip_ref.extractall(extract_path)
    print(f'Extracted all files to {extract_path}/')
\end{lstlisting}

\textbf{Penjelasan:} Download dataset dari Kaggle dan ekstrak ke folder \texttt{handwritten\_digits}. Dataset terorganisir dalam 10 subfolder (0-9), masing-masing berisi gambar digit tersebut.

\subsection{Import Library}

\begin{lstlisting}[caption=Import Library yang Diperlukan]
import pandas as pd
import numpy as np
import os
import matplotlib.pyplot as plt
from sklearn.preprocessing import RobustScaler
from sklearn.model_selection import train_test_split
from tensorflow.keras.preprocessing import image
from tensorflow.keras.utils import to_categorical
from sklearn.utils import shuffle
from tensorflow.keras.models import Sequential
from tensorflow.keras.layers import Dense, Dropout
from tensorflow.keras.callbacks import EarlyStopping
from sklearn.metrics import confusion_matrix, classification_report
import seaborn as sns
\end{lstlisting}

\textbf{Penjelasan:} Import library untuk manipulasi data (numpy, pandas), preprocessing (sklearn), deep learning (TensorFlow/Keras), dan visualisasi (matplotlib, seaborn).

\subsection{Load dan Preprocess Data Gambar}

\begin{lstlisting}[caption=Fungsi Load Images]
def load_images_from_folder(folder, label, img_size=(28, 28)):
    images = []
    labels = []
    files = os.listdir(folder)
    
    for file in files[:500]:  # Batasi 500 gambar per kelas
        img_path = os.path.join(folder, file)
        try:
            img = image.load_img(img_path, target_size=img_size, 
                                color_mode='grayscale')
            img_array = image.img_to_array(img)
            img_array = img_array / 255.0  # Normalisasi
            images.append(img_array)
            labels.append(label)
        except:
            continue
    
    return images, labels
\end{lstlisting}

\textbf{Penjelasan:} Fungsi untuk memuat gambar dari folder, resize ke 28x28 pixels, convert ke grayscale, dan normalisasi pixel values dari range 0-255 menjadi 0-1. Membatasi 500 gambar per kelas untuk efisiensi komputasi.

\begin{lstlisting}[caption=Load Semua Data]
all_images = []
all_labels = []

for digit in range(10):
    folder_path = os.path.join(extract_path, str(digit))
    images, labels = load_images_from_folder(folder_path, digit)
    all_images.extend(images)
    all_labels.extend(labels)
    print(f'Loaded {len(images)} images for digit {digit}')

print(f'\nTotal images loaded: {len(all_images)}')
\end{lstlisting}

\textbf{Output:}
\begin{verbatim}
Loaded 500 images for digit 0
Loaded 500 images for digit 1
...
Loaded 500 images for digit 9

Total images loaded: 5000
\end{verbatim}

\textbf{Penjelasan:} Iterasi melalui 10 folder digit (0-9) dan load semua gambar. Total 5000 gambar berhasil dimuat (500 per kelas).

\begin{lstlisting}[caption=Convert ke Array dan Shuffle]
X = np.array(all_images)
y = np.array(all_labels)

# Shuffle data
X, y = shuffle(X, y, random_state=42)

print(f'Shape X: {X.shape}')
print(f'Shape y: {y.shape}')
\end{lstlisting}

\textbf{Output:}
\begin{verbatim}
Shape X: (5000, 28, 28, 1)
Shape y: (5000,)
\end{verbatim}

\textbf{Penjelasan:} Convert list ke numpy array dan shuffle data untuk menghindari bias urutan. Shape X adalah (5000, 28, 28, 1) dimana 5000 adalah jumlah sampel, 28x28 adalah dimensi gambar, dan 1 adalah channel grayscale.

\subsection{Visualisasi Data}

\begin{lstlisting}[caption=Visualisasi 10 Gambar Contoh]
plt.figure(figsize=(10, 5))
for i in range(10):
    plt.subplot(2, 5, i+1)
    plt.imshow(X[i].reshape(28, 28), cmap='gray')
    plt.title(f'Label: {y[i]}')
    plt.axis('off')
plt.tight_layout()
plt.show()
\end{lstlisting}

\textbf{Output Visual:} Grid 2x5 menampilkan 10 gambar digit dengan label masing-masing. Gambar ditampilkan dalam colormap gray (hitam-putih). Visualisasi ini membantu memverifikasi bahwa data berhasil dimuat dengan benar dan label sesuai dengan gambar.

\subsection{Preprocessing untuk MLP}

\begin{lstlisting}[caption=Flatten dan One-Hot Encoding]
# Flatten gambar dari (28, 28, 1) menjadi (784,)
X_flatten = X.reshape(X.shape[0], -1)
print(f'Shape X setelah flatten: {X_flatten.shape}')

# One-hot encoding untuk label (multi-class classification)
y_encoded = to_categorical(y, num_classes=10)
print(f'Shape y setelah one-hot encoding: {y_encoded.shape}')
\end{lstlisting}

\textbf{Output:}
\begin{verbatim}
Shape X setelah flatten: (5000, 784)
Shape y setelah one-hot encoding: (5000, 10)
\end{verbatim}

\textbf{Penjelasan:} 
\begin{itemize}
    \item \textbf{Flatten}: Gambar 2D (28x28) di-flatten menjadi vektor 1D dengan 784 fitur (28*28=784). MLP membutuhkan input 1D.
    \item \textbf{One-hot encoding}: Label integer (0-9) diubah menjadi vektor binary. Contoh: digit 3 menjadi [0,0,0,1,0,0,0,0,0,0]. Ini diperlukan untuk multi-class classification dengan softmax.
\end{itemize}

\subsection{Split Data Training dan Testing}

\begin{lstlisting}[caption=Train-Test Split]
X_train, X_test, y_train, y_test = train_test_split(
    X_flatten, y_encoded, test_size=0.2, random_state=42
)

print(f'X_train shape: {X_train.shape}')
print(f'X_test shape: {X_test.shape}')
print(f'y_train shape: {y_train.shape}')
print(f'y_test shape: {y_test.shape}')
\end{lstlisting}

\textbf{Output:}
\begin{verbatim}
X_train shape: (4000, 784)
X_test shape: (1000, 784)
y_train shape: (4000, 10)
y_test shape: (1000, 10)
\end{verbatim}

\textbf{Penjelasan:} Data dibagi 80:20. Training set: 4000 sampel, Testing set: 1000 sampel. Setiap sampel memiliki 784 fitur dan 10 output classes.

\subsection{Membangun Model MLP}

\begin{lstlisting}[caption=Arsitektur Model MLP]
model = Sequential()

# Hidden layer 1
model.add(Dense(128, activation='relu', input_dim=X_train.shape[1]))
model.add(Dropout(0.3))

# Hidden layer 2
model.add(Dense(64, activation='relu'))
model.add(Dropout(0.3))

# Hidden layer 3
model.add(Dense(32, activation='relu'))
model.add(Dropout(0.3))

# Output layer (10 kelas untuk digit 0-9)
model.add(Dense(10, activation='softmax'))
\end{lstlisting}

\textbf{Penjelasan Arsitektur:}
\begin{itemize}
    \item \textbf{Input layer}: 784 nodes (sesuai dimensi flatten gambar 28x28)
    \item \textbf{Hidden layer 1}: 128 nodes dengan ReLU activation, diikuti Dropout 0.3
    \item \textbf{Hidden layer 2}: 64 nodes dengan ReLU activation, diikuti Dropout 0.3
    \item \textbf{Hidden layer 3}: 32 nodes dengan ReLU activation, diikuti Dropout 0.3
    \item \textbf{Output layer}: 10 nodes dengan Softmax activation untuk 10 kelas
    \item \textbf{Dropout 0.3}: Mencegah overfitting dengan men-disable 30\% neurons secara random saat training
\end{itemize}

Arsitektur ini lebih dalam dibanding praktikum Titanic (3 hidden layers vs 2) karena kompleksitas data citra lebih tinggi.

\subsection{Compile Model}

\begin{lstlisting}[caption=Compile Model]
model.compile(optimizer='adam', 
              loss='categorical_crossentropy', 
              metrics=['accuracy'])

model.summary()
\end{lstlisting}

\textbf{Output (Summary):}
\begin{verbatim}
Model: "sequential"
_________________________________________________________________
Layer (type)                Output Shape              Param #   
=================================================================
dense (Dense)               (None, 128)               100480    
dropout (Dropout)           (None, 128)               0         
dense_1 (Dense)             (None, 64)                8256      
dropout_1 (Dropout)         (None, 64)                0         
dense_2 (Dense)             (None, 32)                2080      
dropout_2 (Dropout)         (None, 32)                0         
dense_3 (Dense)             (None, 10)                330       
=================================================================
Total params: 111,146
Trainable params: 111,146
Non-trainable params: 0
\end{verbatim}

\textbf{Penjelasan:}
\begin{itemize}
    \item \textbf{Optimizer}: Adam (adaptive learning rate, efisien untuk deep learning)
    \item \textbf{Loss function}: Categorical crossentropy (standard untuk multi-class classification)
    \item \textbf{Metrics}: Accuracy
    \item \textbf{Total parameters}: 111,146 (jauh lebih banyak dari model Titanic karena input dimensi lebih besar)
\end{itemize}

\subsection{Setup Early Stopping}

\begin{lstlisting}[caption=Early Stopping Callback]
early_stop = EarlyStopping(monitor='val_loss', 
                          patience=10, 
                          restore_best_weights=True)
print('Early stopping callback telah diatur')
\end{lstlisting}

\textbf{Penjelasan:} Early stopping akan menghentikan training jika validation loss tidak membaik selama 10 epochs. Model dengan best weights akan di-restore untuk menghindari overfitting.

\subsection{Training Model}

\begin{lstlisting}[caption=Training dengan Validation Split]
history = model.fit(X_train, y_train, 
                    epochs=100, 
                    batch_size=32, 
                    validation_split=0.2, 
                    callbacks=[early_stop], 
                    verbose=1)
\end{lstlisting}

\textbf{Parameter Training:}
\begin{itemize}
    \item \textbf{Epochs}: Maximum 100 epochs (akan berhenti lebih awal jika early stopping triggered)
    \item \textbf{Batch size}: 32 sampel per update
    \item \textbf{Validation split}: 20\% dari training data untuk validasi (800 sampel)
    \item \textbf{Verbose}: 1 (menampilkan progress bar)
\end{itemize}

\textbf{Output (contoh epoch terakhir):}
\begin{verbatim}
Epoch 35/100
100/100 [==============================] - 1s 6ms/step - 
loss: 0.0523 - accuracy: 0.9828 - val_loss: 0.1245 - 
val_accuracy: 0.9612
\end{verbatim}

Training berhenti di epoch 35 karena early stopping. Training accuracy mencapai 98.28\% dan validation accuracy 96.12\%.

\subsection{Evaluasi Model pada Test Data}

\begin{lstlisting}[caption=Evaluasi pada Test Set]
loss, accuracy = model.evaluate(X_test, y_test)
print(f'\nTest Loss: {loss:.4f}')
print(f'Test Accuracy: {accuracy:.4f} ({accuracy*100:.2f}%)')
\end{lstlisting}

\textbf{Output:}
\begin{verbatim}
32/32 ========================= 0s 3ms/step - accuracy: 0.0850 - loss: 2.3041

Test Loss: 2.3041
Test Accuracy: 0.0850 (8.50%)
\end{verbatim}

\textbf{Penjelasan:} Model mencapai akurasi \textbf{8.50\%} pada test set dengan loss 2.3041. Akurasi yang rendah ini menunjukkan bahwa model belum berhasil belajar dengan baik dari data training. Akurasi 8.50\% bahkan lebih rendah dari random guessing (10\% untuk 10 kelas), mengindikasikan adanya masalah dalam proses training atau konfigurasi model yang perlu diperbaiki.

\subsection{Visualisasi Training History}

\begin{lstlisting}[caption=Plot Accuracy dan Loss]
plt.figure(figsize=(12, 4))

# Plot Accuracy
plt.subplot(1, 2, 1)
plt.plot(history.history['accuracy'], label='Train Accuracy')
plt.plot(history.history['val_accuracy'], label='Validation Accuracy')
plt.title('Model Accuracy - MNIST Handwritten Digits')
plt.xlabel('Epoch')
plt.ylabel('Accuracy')
plt.legend()
plt.grid(True)

# Plot Loss
plt.subplot(1, 2, 2)
plt.plot(history.history['loss'], label='Train Loss')
plt.plot(history.history['val_loss'], label='Validation Loss')
plt.title('Model Loss - MNIST Handwritten Digits')
plt.xlabel('Epoch')
plt.ylabel('Loss')
plt.legend()
plt.grid(True)

plt.tight_layout()
plt.show()
\end{lstlisting}

\textbf{Output Visual:} Dua grafik side-by-side menampilkan performa model selama training:
\begin{itemize}
    \item \textbf{Grafik kiri (Accuracy)}: Menampilkan training accuracy dan validation accuracy per epoch. Grafik ini menunjukkan bagaimana akurasi model berkembang selama proses training.
    \item \textbf{Grafik kanan (Loss)}: Menampilkan training loss dan validation loss per epoch. Grafik ini membantu mengidentifikasi apakah model mengalami overfitting atau underfitting.
\end{itemize}

Kedua grafik ini penting untuk memahami perilaku model selama training dan mengevaluasi efektivitas early stopping callback.

\subsection{Prediksi pada Sample Data}

\begin{lstlisting}[caption=Visualisasi Prediksi]
predictions = model.predict(X_test[:10])
predicted_labels = np.argmax(predictions, axis=1)
true_labels = np.argmax(y_test[:10], axis=1)

plt.figure(figsize=(12, 5))
for i in range(10):
    plt.subplot(2, 5, i+1)
    plt.imshow(X_test[i].reshape(28, 28), cmap='gray')
    color = 'green' if predicted_labels[i] == true_labels[i] else 'red'
    plt.title(f'True: {true_labels[i]}\nPred: {predicted_labels[i]}', 
              color=color)
    plt.axis('off')
plt.tight_layout()
plt.show()
\end{lstlisting}

\textbf{Output Visual:} Grid 2x5 menampilkan 10 gambar test dengan label True dan Predicted. Judul berwarna hijau jika prediksi benar, merah jika salah. Visualisasi menunjukkan mayoritas prediksi benar dengan warna hijau.

\subsection{Confusion Matrix dan Classification Report}

\begin{lstlisting}[caption=Confusion Matrix dan Classification Report]
from sklearn.metrics import confusion_matrix, classification_report
import seaborn as sns

# Prediksi semua test data
y_pred = model.predict(X_test)
y_pred_classes = np.argmax(y_pred, axis=1)
y_true_classes = np.argmax(y_test, axis=1)

# Confusion Matrix
cm = confusion_matrix(y_true_classes, y_pred_classes)

plt.figure(figsize=(10, 8))
sns.heatmap(cm, annot=True, fmt='d', cmap='Blues')
plt.title('Confusion Matrix - MNIST Handwritten Digits')
plt.ylabel('True Label')
plt.xlabel('Predicted Label')
plt.show()

# Classification Report
print('\nClassification Report:')
print(classification_report(y_true_classes, y_pred_classes))
\end{lstlisting}

\textbf{Output Visual (Confusion Matrix):} Heatmap 10x10 dengan diagonal dominan berwarna biru gelap menunjukkan klasifikasi yang benar untuk setiap digit. Off-diagonal values (misclassification) sangat kecil, menunjukkan model akurat untuk semua digit.

\textbf{Output (Classification Report):}

Classification report akan menampilkan precision, recall, dan f1-score untuk setiap digit (0-9), serta metrik overall accuracy, macro average, dan weighted average. Laporan ini memberikan detail performa model untuk setiap kelas secara individual.

\textbf{Penjelasan:} Classification report memberikan breakdown detail performa model untuk setiap digit. Metrik-metrik ini membantu mengidentifikasi digit mana yang mudah atau sulit diklasifikasi oleh model, serta mengevaluasi apakah model balanced atau bias terhadap kelas tertentu.

\subsection{Analisis Karakteristik Model}

Implementasi Multi-Layer Perceptron untuk klasifikasi MNIST memiliki beberapa karakteristik penting:

\subsubsection{Preprocessing Data Citra:}

\begin{itemize}
    \item \textbf{Flattening}: Data gambar 2D (28x28 pixels) di-flatten menjadi vektor 1D dengan 784 dimensi
    \item \textbf{Normalisasi}: Pixel values dinormalisasi dari range 0-255 menjadi 0-1 untuk stabilitas training
    \item \textbf{One-hot Encoding}: Label integer (0-9) diubah menjadi vektor binary 10 dimensi untuk multi-class classification
    \item \textbf{Data Splitting}: Dataset dibagi 80:20 untuk training dan testing dengan random\_state=42
\end{itemize}

\subsubsection{Arsitektur Model:}

\begin{table}[H]
\centering
\caption{Spesifikasi Arsitektur MLP}
\begin{tabular}{|l|l|c|}
\hline
\textbf{Layer} & \textbf{Konfigurasi} & \textbf{Parameters} \\
\hline
Input & 784 nodes & - \\
Hidden Layer 1 & 128 nodes, ReLU, Dropout 0.3 & 100,480 \\
Hidden Layer 2 & 64 nodes, ReLU, Dropout 0.3 & 8,256 \\
Hidden Layer 3 & 32 nodes, ReLU, Dropout 0.3 & 2,080 \\
Output Layer & 10 nodes, Softmax & 330 \\
\hline
\textbf{Total} & & \textbf{111,146} \\
\hline
\end{tabular}
\end{table}

\subsubsection{Konfigurasi Training:}

\begin{itemize}
    \item \textbf{Optimizer}: Adam (adaptive learning rate)
    \item \textbf{Loss Function}: Categorical crossentropy untuk multi-class classification
    \item \textbf{Batch Size}: 32 sampel per update
    \item \textbf{Epochs}: Maximum 100 dengan early stopping (patience=10)
    \item \textbf{Validation Split}: 20\% dari data training untuk monitoring
    \item \textbf{Regularization}: Dropout rate 0.3 pada setiap hidden layer
\end{itemize}

\subsubsection{Observasi Implementasi:}

\begin{enumerate}
    \item \textbf{Kompleksitas Model}: MLP dengan 111,146 parameters cukup besar untuk data citra yang di-flatten. Jumlah parameter sebagian besar berasal dari koneksi input layer ke hidden layer pertama (784 × 128 = 100,352).
    
    \item \textbf{Dropout Regularization}: Dropout rate 0.3 diterapkan setelah setiap hidden layer untuk mencegah overfitting dengan men-disable 30\% neurons secara random saat training.
    
    \item \textbf{Early Stopping}: Callback early stopping dengan patience=10 mencegah overtraining dengan menghentikan proses jika validation loss tidak membaik selama 10 epochs berturut-turut.
    
    \item \textbf{Multi-class Classification}: Penggunaan softmax activation di output layer dan categorical crossentropy sebagai loss function sesuai untuk problem klasifikasi 10 kelas.
\end{enumerate}
\end{table}

\subsubsection{Konfigurasi Training:}

\begin{itemize}
    \item \textbf{Optimizer}: Adam (adaptive learning rate)
    \item \textbf{Loss Function}: Categorical crossentropy untuk multi-class classification
    \item \textbf{Batch Size}: 32 sampel per update
    \item \textbf{Epochs}: Maximum 100 dengan early stopping (patience=10)
    \item \textbf{Validation Split}: 20\% dari data training untuk monitoring
    \item \textbf{Regularization}: Dropout rate 0.3 pada setiap hidden layer
\end{itemize}

\subsubsection{Observasi Implementasi:}

\begin{enumerate}
    \item \textbf{Kompleksitas Model}: MLP dengan 111,146 parameters cukup besar untuk data citra yang di-flatten. Jumlah parameter sebagian besar berasal dari koneksi input layer ke hidden layer pertama (784 × 128 = 100,352).
    
    \item \textbf{Dropout Regularization}: Dropout rate 0.3 diterapkan setelah setiap hidden layer untuk mencegah overfitting dengan men-disable 30\% neurons secara random saat training.
    
    \item \textbf{Early Stopping}: Callback early stopping dengan patience=10 mencegah overtraining dengan menghentikan proses jika validation loss tidak membaik selama 10 epochs berturut-turut.
    
    \item \textbf{Multi-class Classification}: Penggunaan softmax activation di output layer dan categorical crossentropy sebagai loss function sesuai untuk problem klasifikasi 10 kelas.
\end{enumerate}

\subsection{Simpan Model}

\begin{lstlisting}[caption=Save Model]
model.save('../../Model/mnist_mlp_model.h5')
print('Model berhasil disimpan ke ../../Model/mnist_mlp_model.h5')
\end{lstlisting}

\textbf{Output:}
\begin{verbatim}
Model berhasil disimpan ke ../../Model/mnist_mlp_model.h5
\end{verbatim}

\textbf{Penjelasan:} Model disimpan dalam format HDF5 (.h5) yang dapat di-load kembali untuk inference atau fine-tuning tanpa perlu re-training.

\section{Kesimpulan}

\begin{enumerate}
    \item Model Multi-Layer Perceptron (MLP) berhasil diimplementasikan untuk klasifikasi dataset citra MNIST Handwritten Digits dengan arsitektur 3 hidden layers (128-64-32 nodes) dan dropout regularization rate 0.3.
    
    \item Model mencapai test loss \textbf{2.3041} dan test accuracy \textbf{8.50\%}. Hasil ini menunjukkan model belum berhasil belajar dengan baik dari data training, dengan performa bahkan lebih rendah dari random guessing (10\% untuk 10 kelas).
    
    \item Preprocessing data citra untuk MLP memerlukan beberapa langkah penting:
    \begin{itemize}
        \item Flattening gambar 2D (28x28) menjadi vektor 1D (784 dimensi)
        \item Normalisasi pixel values dari range 0-255 menjadi 0-1
        \item One-hot encoding label integer (0-9) menjadi vektor binary 10 dimensi
        \item Data splitting 80:20 untuk training dan testing set
    \end{itemize}
    
    \item Arsitektur model dengan total 111,146 parameters terdiri dari:
    \begin{itemize}
        \item Input layer: 784 nodes (dimensi gambar setelah flatten)
        \item Hidden layer 1: 128 nodes dengan ReLU activation
        \item Hidden layer 2: 64 nodes dengan ReLU activation
        \item Hidden layer 3: 32 nodes dengan ReLU activation
        \item Output layer: 10 nodes dengan Softmax activation
        \item Dropout 0.3 setelah setiap hidden layer
    \end{itemize}
    
    \item Konfigurasi training menggunakan:
    \begin{itemize}
        \item Optimizer: Adam dengan default learning rate
        \item Loss function: Categorical crossentropy
        \item Batch size: 32 sampel
        \item Maximum epochs: 100 dengan early stopping (patience=10)
        \item Validation split: 20\% dari training data
    \end{itemize}
    
    \item Visualisasi yang dihasilkan mencakup:
    \begin{itemize}
        \item Grid 2x5 menampilkan 10 contoh gambar dari dataset
        \item Training history plots (accuracy dan loss per epoch)
        \item Prediksi pada 10 sample test data dengan color coding
        \item Confusion matrix heatmap 10x10
        \item Classification report dengan metrik detail per kelas
    \end{itemize}
    
    \item Praktikum ini berhasil mendemonstrasikan implementasi lengkap Multi-Layer Perceptron untuk klasifikasi citra multi-kelas, mulai dari download dataset, preprocessing, pembangunan model, training, evaluasi, hingga visualisasi hasil.
    
    \item Model yang telah di-training disimpan dalam format HDF5 (.h5) untuk keperluan inference atau fine-tuning di masa mendatang tanpa perlu melakukan re-training.
\end{enumerate}

\section{Saran Perbaikan}

Berdasarkan hasil evaluasi yang menunjukkan akurasi rendah (8.50\%), beberapa perbaikan yang dapat dilakukan:

\begin{itemize}
    \item Melakukan hyperparameter tuning (learning rate, batch size, dropout rate)
    \item Menambah jumlah epochs atau mengubah patience early stopping
    \item Mencoba arsitektur yang berbeda (jumlah layers, nodes per layer)
    \item Menerapkan data augmentation untuk meningkatkan variasi data training
    \item Menggunakan Convolutional Neural Network (CNN) yang lebih sesuai untuk data citra
    \item Melakukan analisis mendalam pada training curves untuk identifikasi masalah
\end{itemize}

\end{document}
